\section{Transfer of the surface-layer closure to the Fortran reference}
\label{app:fortran-slab}

The closure of Sect.~\ref{sec:structural-slab} acts on physical,
grid-computed quantities only, the surface sensible heat flux and the
air temperature entering the surface exchange, and is therefore not
tied to the differentiable implementation. This appendix documents its
transfer to the Fortran offline driver of the reference code base and the
bitwise invariance of the inactive closure and its transfer.

\subsection{Implementation in the offline driver}

The recurrence is added to the driver level (\texttt{cpg1s.F90}), outside the
SURF physics itself, in four changes: (i)~a persistent anomaly array
per grid point; (ii)~immediately before each call to the SURF physics, the
lowest-level air temperature argument at the incoming time level is
incremented by the anomaly, and the increment is removed again after the
call; (iii)~the anomaly is advanced with the
identical discretization as Eq.~(\ref{eq:slab}),
$\Delta T_{n+1} = \beta_n \Delta T_n + (1-\beta_n)(-\lambda_s H_n)$, with the
shortwave-blended per-step relaxation factor $\beta_n$ evaluated from the
driver's own surface downward shortwave argument, and $H_n$
the grid-mean surface sensible heat flux of the step just completed
(tile-fraction weighted, downward positive, the same quantity the torch
model differentiates through, exported from the vertical-diffusion
interface of the reference physics); and (iv)~environment-variable switches
with compiled-in defaults $\lambda_s = 0$ and the tied timescale, so the
unmodified driver behaviour is the compiled-in default. The constants are
shared with the torch implementation
($\rho c_p = 1.2\times1005.7$~J\,m$^{-3}$\,K$^{-1}$,
$S_{ref} = 50$~W\,m$^{-2}$, $\Delta t = 1800$~s), and the update ordering
matches: the anomaly applied at step $n{+}1$ is built from the flux and
shortwave of step $n$.

\subsection{Bitwise invariance of the inactive closure}

With $\lambda_s = 0$ the patch must be numerically invisible. A one-day
integration of the patched executable (48 steps, 97\,709 columns) was
compared record by record against an archive produced by the unpatched
executable from the same restart and forcing: all 24 output variables of
the 3-hourly history file agree bitwise, and the run log confirms the
closure inactive ($\lambda_s = 0$ reported at initialization).

\subsection{Matched-protocol evaluation}

All Fortran evaluations use year-long integrations started from the common
spun-up restart of Sect.~\ref{sec:modis-evaluation}; the calibrated
constants of the torch configuration ($\lambda_s =
0.05$~K\,(W\,m$^{-2}$)$^{-1}$, $\tau_{\mathrm{day}} = 6$~h,
$\tau_{\mathrm{night}} = 3$~h) enter the driver through its environment
switches without re-tuning. Because the driver history is 3-hourly, the
Fortran runs and a re-sampled torch evaluation sample the nearest
3-hourly record at each MODIS view time; the observation operator is the
driver's grid-mean skin temperature (\texttt{AvgSurfT}), which is by
construction the same tile-fraction-weighted mean as the torch
\emph{radiative skin temperature}. Under this matched protocol the
year-long unpatched Fortran integration and the torch integration of the
prior configuration agree to within $0.001$~K in bias and RMSE in all four
window--pass combinations (e.g.\ held-out day bias $-8.068$ versus
$-8.068$~K, RMSE $11.217$ versus $11.218$~K), so the MODIS error
statistics of the two implementations are indistinguishable at this
protocol before the closure is activated. What is verified to date is the
numerical fidelity of the transfer (bitwise invariance of the inactive
closure) and the indistinguishability of the prior statistics; full-year
Fortran statistics of the calibrated configuration ($\lambda_s = 0.05$,
$\tau_{\mathrm{day}} = 6$~h, $\tau_{\mathrm{night}} = 3$~h) are not yet
available and are not anticipated here. All conclusions of
Sect.~\ref{sec:structural-slab} rest on the torch implementation.
